\documentclass[11pt,letterpaper]{article}
\usepackage[T1]{fontenc}
\usepackage[utf8]{inputenc}
\usepackage{lmodern}
\usepackage[margin=0.86in]{geometry}
\usepackage{microtype}
\usepackage{booktabs,longtable,tabularx,array}
\usepackage{enumitem}
\usepackage{xcolor}
\usepackage{hyperref}
\usepackage{xurl}
\definecolor{DocBlue}{HTML}{16435F}
\hypersetup{colorlinks=true,urlcolor=DocBlue,linkcolor=DocBlue}
\setlength{\parindent}{0pt}
\setlength{\parskip}{0.54em}
\setlist[itemize]{leftmargin=*,itemsep=3pt}
\begin{document}
\begin{center}
{\LARGE\bfseries Enterprise AppSec Governance Architecture Description}\par
\vspace{5pt}{\large ISO/IEC/IEEE 42010-aligned proposed architecture}\par
\vspace{9pt}{\small Engineering reference | Proposed design | Not a deployment or approval record}\par
\end{center}
\tableofcontents
\clearpage
\section{Entity of interest and boundary}
Entity: enterprise security triage decision capability consuming existing read-only Architecture Discovery evidence. Its scope excludes the current GitHub runtime itself and does not convey merge authority.


\section{Stakeholders and concerns}
\begin{itemize}
\item Developer: unblocking legitimate changes with defensible evidence.
\item AppSec/security manager: risk acceptance authority, false-positive correctness and traceable approvals.
\item Platform owner: least privilege, token boundaries and deployment safety.
\item Auditor: evidence provenance, deterministic dispositions, expiry and review history.
\end{itemize}

\section{Architecture views}
\begin{itemize}
\item Context: engineering repositories, read-only collector, policy evaluation service, GitHub security API and independent reviewer.
\item Logical: inventory adapter, alert adapter, policy bundle, exception store, audit ledger, future privileged execution boundary.
\item Information: repository/PR/alert/exception/policy/evidence identifiers with E1/E2/E3 provenance.
\item Security: untrusted PR data isolated from privileged authorization and immutable policy.
\item Operations: manual escalation, shadow mode, error processing, reconcile and rollback.
\end{itemize}

\section{Correspondences}
Every alert used in a decision must identify its repository, rule and analysis. Every exception must correspond to a separate authorized decision and its scope. Every ALLOW\_MERGE classification must reference completed scans while explicitly deferring to GitHub's effective rules and approvals.


\section{Open questions and decisions}
Policy language, service hosting, storage residency and event triggers remain proposed ADR topics. Deployment, ownership and controls are not approved merely because this AD exists.


\section{Verification}
Compare design views with actual GitHub permissions and current API contracts, run synthetic negative tests, confirm E1/E2/E3 state distinction and test that read-only collectors cannot issue write operations.


\section{Integrated requirements and evidence correspondence}
The documentation sufficiency overlay in the security companion suite maintains stable requirement IDs, design coverage, evidence ownership, verification IDs and L1/L2/L3 readiness outcomes. For each architecture relationship, preserve source commit, collector and extraction version, UTC observation, stable repository/component identity, and E1, E2 or E3 semantics. An inferred owner or component relationship is never a false-positive validation.

Each derived AppSec disposition must additionally trace the alert instance, PR revisions, configured scan completeness, distinct required checks, effective ruleset, reviews and queue obligations, immutable policy identity and, where applicable, an independently verified exception. The output of an offline evaluator remains advisory, with GitHub the authority for merge status.

\section{Reviewability, constraints and quality attributes}
Quality concerns are completeness, freshness, resilience to limited API access, reproducibility, confidentiality and separation of duties. A scope-limited collector cannot approve its own security exceptions. An executor capable of alert mutation is outside the deployed system boundary until separately authorized, secured and tested.

Documentation sign-off and acceptance of proposed ADRs are separate from pilot entry and production go-live. Verification evidence must distinguish observed positive/negative tests from drafts, synthetic fixtures and empty forms. Actual enterprise instance permissions, edition support, records retention policy, service SLOs, recovery and named reviewers are outstanding evidence obligations.

\section*{Document status and provenance}
This document is a proposed reference design. Actual GitHub feature availability, organization policies, approvals, and runtime performance must be independently verified before adoption. The architecture inventory is not proof of vulnerability exploitability or correctness of a dismissal.
\end{document}
